\chapter{Conclusion}
\label{chap:11-conclusion}

\section{The answer}
\label{sec:reponse}

\begin{keybox}
Joelle Pineau was heard on what accompanies the result: the code link, the declared compute, the
effect of a setting and, on weaker evidence, the error bar. The bar has spread without uncertainty computed over random seeds
following it detectably. The tuning of baselines shows no sign of progress.
\end{keybox}

This answer applies to classical RL, read by AI agents in fifty papers on each side, then compared
with our own reading of 2018 using the same grids
(\cref{chap:03-sondage-methode,chap:04-sondage-resultats}). It also rests on experiments rerun with
the algorithms of 2026
(\cref{chap:05-experiences-methode,chap:06-experiences-resultats,chap:07-modeles-de-langue}). The
numbers it rests on come from those chapters. This conclusion says what they establish and refers
back to them.

\section{What the measurements establish}
\label{sec:nombres}

\paragraph{Four practices have changed} The link to the code, the compute stated in figures, the
sensitivity to a setting and reported uncertainty pass the Holm correction in the main comparison
(\cref{tab:sondage-pineau,tab:sondage-extension}). Code and compute survive every cut we tried,
sensitivity nearly every one. Code keeps its verdict when both years are read on the same kind of
version, proceedings or arXiv (\cref{sec:sr-robustesse}). Reported uncertainty is the most fragile
change: it no longer passes the correction once the two papers of the comparison read partly by
another model are removed. Its verdict therefore depends on two papers. Since the readers of 2018 are
not those of 2026, eight fresh readers reread both years: the rises in code, compute and sensitivity
hold, while the rise in reported uncertainty keeps its direction without being established
(\cref{sec:sr-annees}). A reader from another model family finds the direction of all four rises.
When its \enquote{yes} answers without a quotation found in the paper count as \enquote{no}, it does
not confirm the rise in reported uncertainty (\cref{sec:sr-gemma}). Criterion X8 shows what that last rise
measures. The share of papers that compute their uncertainty
over random seeds goes from 35\% to 42\%, a gap the survey does not detect: the error bar has spread
faster than the bar computed over independent training runs (\cref{sec:sr-principale}). Read on 20
papers per year between the two dates, the four practices rise slowly, with no jump right after the
checklist that NeurIPS has required since 2019 (\cref{tab:trajectoire}).

\paragraph{Judging the result has barely moved} The tuning of baselines stagnates: the hyperparameters given for the
baselines (Q4), the same tuning care announced (X4) and the origin of their code (Y5) do not rise.
The version of the benchmarks (Y7), the definition of the performance measure (Y1) and the selection
of the best runs (X3) do not move either. The statistical test, held-out evaluation, the
hyperparameters of the method and the way they are chosen rise by 6 to 22 points depending on the
criterion, with none of these rises passing the correction. The survey could detect with 80\% power
only a rise of 18 to 30 points depending on the criterion, 26 to 39 after correction
(\cref{sec:val-stat}). It does not say that these practices stagnate. It says
that it cannot know.

\paragraph{A ticked box does not guarantee the practice} The 18 NeurIPS 2025 papers that answer yes to their
checklist question on statistical significance apply no test. On code and on compute, some of the
\enquote{yes} answers have neither a link nor a figure behind them (\cref{sec:sr-declarations}).

\paragraph{Chance fools the eye, sorting fools the test} Two groups of five runs drawn at random from those of one
program rarely fool a statistical test. Sorted, the five best against the five worst fool it from
26\% to nearly 100\% of the time depending on the (program, task) pair, 31\% in the main
case, PPO on \task{HalfCheetah} (\cref{sec:e1}). Keeping the best three of ten runs inflates the
mean by 32\% at the median over the sixteen pairs (\cref{tab:e4}). The $t$ of 9.09 from 2018 probably
comes from the script published by Peter Henderson, which measures the spread of the episodes and
ignores the spread of the random seeds. The published ($t$, $p$) pair excludes the other readings of
the computation more than it points to this script: the competing reading, the mean of the $p$
values, has no well-defined degrees of freedom. Applied to our TRPO runs, the script exceeds 9 in
38\% of random draws (\cref{tab:b1}).

\paragraph{The best program depends on the task, the budget and the physics} Across our four programs, the best one changes with the task.
Yet PPO has five times more steps than SAC and CrossQ. At equal budget and with twenty random seeds,
CrossQ leads on all four tasks, but the order of two pairs still reverses from one task to another in
a settled way. The reversal of PPO JAX and SAC even survives the Bonferroni correction
(\cref{sec:e2}). The two implementations of PPO
differ in a settled way only on \task{Swimmer}, by a gap that is small in score points
(\cref{sec:e3}). Under a physics modified by 20\%, CrossQ beats each of the three other programs on
\task{Hopper}, at ten as at twenty random seeds. PPO JAX loses there against SAC, which it beat in
the training physics (\cref{sec:e7}).

\paragraph{A published recipe has aged} The agents trained with the recipe published by RL
Baselines3 Zoo stay upright on \task{Hopper} without moving forward. The agent that the Zoo itself
published in 2022, with the same settings and an older software stack, moves forward and scores
close to our default agents. The recipe worked in its original stack and has aged (\cref{sec:e6}).

\paragraph{On language models, no GRPO gain is established} With five random seeds, no gain within a configuration can be
established, since the exact sign test does not go below 0.0625. One contrast is established: for
Llama, the true reward does better than the randomly drawn reward, because the latter destroys the
model. Under the true reward, Llama loses greedy accuracy and pass@8 in each of the five random seeds,
while its pass@1 rises slightly. For Qwen, the pass@1 gain is explained by the end of truncated
answers: the accuracy of finished answers does not rise. Qwen also ran with a repetition penalty that
Llama did not have (\cref{chap:07-modeles-de-langue}). A rerun of ten random seeds per configuration
waits for the GPU (\cref{sec:val-en-cours}).

\section{Recommendations}
\label{sec:recommandations}

The following recommendations follow from the measurements of this report. Each names the flaw it
targets and the result that justifies it.

\subsection{For authors}

\begin{enumerate}
  \item \textbf{State the number of random seeds and run at least ten.} Of 41 experimental RL papers
        in the first 2026 draw, 20 do not say how many runs they launch. The median of the others is
        five. With five random seeds, a gap of less than 20\% reads backwards nearly one time in
        three, more than one time in three when twenty random seeds set the true order or when the
        order is judged out of sample (E5, flaws F9 and F10).
  \item \textbf{Never sort or select runs.} Report every run, including those that fail. State which
        were excluded (E1, E4, criteria X3 and Y6).
  \item \textbf{Test across random seeds.} The independent unit is the
        complete training run. A test that counts episodes declares a difference between two groups
        of the same code (B1, flaw F11). The same rule holds for the error bar: in 2026, fewer than
        one experimental paper in two computes it over random seeds (X8).
  \item \textbf{Show an interval next to every curve and say what it is computed over.} Name the
        standard deviation, the standard error or the confidence interval, then the spread it
        summarises: 14 of the 44 papers of 2026 that draw a bar do not say what it is computed over
        (criteria X2, X8, Y2, Y11).
  \item \textbf{Compare at equal budget, or state the difference in budget.} A ranking at unequal
        budgets also compares the budgets (E2 at equal budget).
  \item \textbf{Correct for multiple comparisons and say what the sample could see.} Four tasks
        tested without correction made E3 read two differences, only one of which survives the Holm
        correction. A \enquote{no difference} is read together with the smallest difference the
        sample detects (E3, \cref{sec:val-stat}).
  \item \textbf{Tune the baselines with the same care as the method and say so.} Give the tuning
        budget of each baseline, its final hyperparameters and the version of the libraries in which
        its recipe was tuned (E6, criteria Q4, X4, Y5).
  \item \textbf{Evaluate under held-out conditions.} Environment random seeds never seen, modified
        dynamics or a distinct test distribution (E7, criteria Q7, Y8).
  \item \textbf{Give the versions.} The version of each benchmark and of each dependency, in a frozen
        environment file (criteria Y7, Y10).
  \item \textbf{For a language model: several model families, the full decoding setup and several
        measures.} A result obtained on one family informs about that family. A setting inherited
        from the model's configuration file, such as a repetition penalty, changes what is sampled. A
        gain in accuracy is read together with the length limit, the truncated share before and
        after, the accuracy of the complete answers, the greedy answer and pass@$k$
        (\cref{chap:07-modeles-de-langue}, criterion X6).
  \item \textbf{Write the protocol before the first run and publish every change.} For each change,
        state what had already been seen when it was decided.
  \item \textbf{Recompute every published number from its files.} A script that checks every number
        in the text against its result file costs little and catches drift
        (\cref{sec:verification}).
\end{enumerate}

\subsection{For reviewers}

The seven questions of the blog article are meant for any reader. A reviewer can ask them of every
submission.
\begin{enumerate}
  \item How many random seeds? With fewer than five runs, or with \enquote{the best runs}, the result
        says little.
  \item What uncertainty? A mean with no interval and no band around the curve does not allow a gap
        to be judged.
  \item Which baseline, tuned how? If the method got a month of tuning and the baseline its default
        values, the gap measures the month of tuning.
  \item What budget? Comparing two methods trained on different numbers of steps also compares those
        numbers of steps.
  \item Evaluated where? An agent judged on the situations of its training may have learned them by
        heart.
  \item What code? Without published code, the reader does not know what the paper really did.
  \item What compute? A result that nobody else can afford to recompute cannot be reproduced.
\end{enumerate}
The reviewer can add an eighth question, which this report justifies: over which unit are the
uncertainty and the test computed? A bar or a test over episodes, or over the stacked points of a
curve, does not measure the variability between training runs.

\subsection{For conferences}

\begin{enumerate}
  \item \textbf{Ask for the test as well as the bars.} The NeurIPS question on statistical
        significance accepts error bars \enquote{or other appropriate information}. The 18 papers of
        our sample that answer yes to it apply no test. A question that asks for the unit of the test and the number
        of independent training runs would demand more.
  \item \textbf{Check a ticked box on a sample.} Some papers answer yes on code with no link, on
        compute with no figure. A check drawn at random on a share of the accepted papers would
        measure the gap between declaration and practice.
  \item \textbf{Require the link to the code in the camera-ready version.} \enquote{Code will be
        released} is no substitute for a link. A link added only to the arXiv version escapes the
        proceedings.
  \item \textbf{Ask for the version of each benchmark.} This practice has not moved detectably since
        2018 (23\% then 26\%), although the three agentic AI models that rated the cost of each
        practice (\cref{chap:04-sondage-resultats}) judge it nearly free.
  \item \textbf{Support reproducibility challenges.} The ICLR 2018 challenge presented by Pineau
        showed that reproduction by third parties is feasible: among the participants who answered
        the survey, 54.5\% reproduced part of a paper, 32.7\% most of it
        \parencite[slides 37 to 43]{pineau2018iclr}. Running it again by broad families of results,
        with a protocol written in advance, would measure what the checklist does not.
\end{enumerate}

\section{Future work}
\label{sec:futurs}

\begin{itemize}
  \item \textbf{ICLR in the trajectory.} The years 2019 to 2024 lack ICLR. A reading of its
        proceedings, through a download route that OpenReview accepts, would complete the series.
  \item \textbf{A larger survey on the criteria that rise without being established.} The
        statistical test, held-out evaluation and the method for choosing settings rise by 6 to 22
        points without the survey being able to establish it. A sample large enough to detect a rise of ten
        points would settle it.
  \item \textbf{The papers that apply what they call RL to large language models, read with the same
        grid} on a sample of comparable size, would tell whether held-out evaluation, which is
        progressing there, comes with the other practices.
  \item \textbf{The blind human rereading.} Thirty papers, fifteen per year, await our
        blind reading. The script that will compare it with the agents is written and refuses
        to run until all thirty sheets are filled.
  \item \textbf{A new complete training run on another platform}, Linux on x86-64, would test the
        statistics of our experiments with other draws (\cref{sec:reentrainer}).
  \item \textbf{Other families of tasks} than MuJoCo locomotion: image-based games, offline RL, real
        robotics, where the conclusions of E2 to E7 have yet to be tested.
  \item \textbf{The GRPO strand at a larger scale.} If the rerun with ten random seeds establishes
        its contrast, the same protocol on 7-billion-parameter models would bring the strand closer to
        \textcite{shao2025spurious}.
\end{itemize}

\paragraph{What we corrected}
In 2018, Pineau asked researchers to be as critical of their good results as of their bad ones. We
looked for our own errors and corrected each one. In E6, our code gave the reward normalisation of
the Zoo a discount factor of 0.99 instead of 0.999: the corrected recipe does no better in any
detectable way (\cref{sec:e6}). The E2 count counted triplets instead of the pairs that reverse: the
strict count finds three pairs instead of fourteen changes and the verdict holds (\cref{sec:e2}). The
factor that links the $t$ of Henderson's script to that of a test over random seeds is $\sqrt2/\rho$
for our ten episodes, not $1/\rho$ (\cref{sec:b1}). In the GRPO strand, the Qwen gain, clear at three
random seeds, is no longer clear at five: it comes from the length limit, which cut its answers off.
Qwen also decoded with a repetition penalty that Llama did not have and the pilot that set the task
had read the test set: the published figures stay the same, while the comparison between families
mixes the family with the decoding (\cref{chap:07-modeles-de-langue}). The repository keeps a record
of every correction for anyone who wants to redo the work.

\section{Opening}\label{sec:ouverture}
As early as 2019, Michael I. Jordan described artificial intelligence as an engineering discipline
still to be built, rather than a revolution already made \parencite{mijordan2019revolution,mijordan2019drai}.
He returns to that idea about a recent incident in which AI agents were left free to act: they are,
he says, \enquote{mostly bad engineering practice}. He compares: \enquote{in chemical engineering,
we would never have had a situation where explosions were allowed to happen to see how far things
could go wrong}
\parencite[our translation from the French]{mijordan2026liberation}.

Pineau's three Rs state the same demand at the scale of a paper. A reproducible result, reusable
parts, a robust conclusion: that is what an engineer expects of any structure before using it. Our
survey shows that RL has adopted the outward marks of that discipline, the published code, the
declared compute, the error bar. The way a result is judged has changed much less: the number of
runs, the statistical test, the held-out evaluation.

Michael I. Jordan also sees AI as an economic and social phenomenon, made of producers, consumers and
people who decide in which domains they accept it. For AI programs that learn by trial and error,
their trust will come through results that others can redo. The methods exist. A protocol written in
advance costs little next to a false conclusion.
